%! Piecewise \& Step Functions — Unit 2, Lesson 2.4 — Warm-Up
\documentclass[10pt]{article}
\usepackage{algebra2-article}
\usepackage{algebra2-boxes}
\newcommand{\TallMath}[1]{$\displaystyle #1\rule[-1.4em]{0pt}{3.2em}$}

\begin{document}

\pageheader{Unit 2, Lesson 2.4}{Warm-Up}
\namedateperiod

\vspace{0.10in}

\begin{notesbox}{Getting Ready: Rules, Two Pieces, and Endpoints}
\small Today a function can follow more than one rule. These three quick reviews are the tools you'll
need.

\vspace{4pt}
\textbf{1. Evaluate the rule.} For $g(x)=2x-1$, find each output.
\begin{center}
\renewcommand{\arraystretch}{1.25}
\begin{tabular}{|c|c|c|c|}
\hline
$x$ & $-2$ & $0$ & $3$ \\ \hline
$g(x)=2x-1$ & \blank{0.9cm} & \blank{0.9cm} & \blank{0.9cm} \\ \hline
\end{tabular}
\end{center}
And for $h(x)=|x|$: \quad $h(-4)=\blank{1.0cm}$, \quad $h(0)=\blank{1.0cm}$.

\vspace{6pt}
\textbf{2. The V is two lines (from Lesson 2.3).} The parent $y=|x|$ is made of two straight lines
pasted at the vertex. Fill in each rule.
\[ y=|x| \text{ behaves like } y=\blank{1.2cm} \text{ when } x\ge 0, \text{ and like }
   y=\blank{1.2cm} \text{ when } x<0. \]

\vspace{4pt}
\textbf{3. Endpoints and inequalities.} A number line uses a \emph{closed} dot $\bullet$ to
\emph{include} a point and an \emph{open} dot $\circ$ to \emph{exclude} it.
\begin{itemize}[leftmargin=1.6em, itemsep=3pt]
  \item Is $x=1$ a solution of $x<1$? \blank{1.4cm} \quad Is $x=1$ a solution of $x\le 1$? \blank{1.4cm}
  \item To show ``$x\ge 2$'' on a number line, the dot at $2$ should be \blank{2.4cm} (open / closed).
\end{itemize}
\end{notesbox}

\end{document}
